\documentclass[11pt,reqno]{amsart}
\usepackage{lmodern}
\usepackage[T1]{fontenc}
\usepackage[utf8]{inputenc}
\usepackage{amssymb,amsmath,amsthm,mathtools}
\usepackage{booktabs}
\usepackage{longtable}
\usepackage{array}
\usepackage{float}
\usepackage{enumitem}
\usepackage{microtype}
\usepackage{xcolor}
\usepackage[margin=1in]{geometry}
\usepackage[colorlinks=true,linkcolor=blue!50!black,citecolor=blue!50!black,urlcolor=blue!50!black]{hyperref}

\theoremstyle{plain}
\newtheorem{theorem}{Theorem}
\newtheorem{thm}[theorem]{Theorem}
\newtheorem{proposition}{Proposition}
\newtheorem{lemma}{Lemma}
\newtheorem{corollary}{Corollary}
\theoremstyle{definition}
\newtheorem{definition}{Definition}
\newtheorem{example}{Example}
\theoremstyle{remark}
\newtheorem{remark}{Remark}
\newtheorem{observation}{Observation}
\newtheorem{conjecture}{Conjecture}
\newtheorem{question}{Question}

\newcommand{\Z}{\mathbb{Z}}
\newcommand{\F}{\mathbb{F}}
\newcommand{\ord}{\operatorname{ord}}
\DeclareMathOperator{\lcm}{lcm}
\newcommand{\gen}[1]{\langle #1 \rangle}
\newcommand{\dvd}{\mid}
\newcommand{\Sa}{S_a}
\newcommand{\half}{\tfrac12}
\DeclareMathOperator{\popcount}{w}

\title[On the congruences $2^{n}\equiv-1$ and $2^{n}\equiv-3\pmod{n+1}$]{On the congruences $2^{n}\equiv-1$ and $2^{n}\equiv-3 \pmod{n+1}$}
\author{Robin Rovenszky}
\author{notxxdog}
\date{}
\subjclass[2020]{Primary 11A07, 11Y11; Secondary 11A51, 11N25, 11Y05}
\keywords{Power-of-two congruences, Fermat pseudoprimes, multiplicative order, discrete logarithm, admissibility sieve, Wieferich primes, OEIS A245728}
\thanks{Prepared with substantial assistance from an artificial-intelligence system; see the appendix on the use of artificial intelligence. Author contacts: Robin Rovenszky (\texttt{robincodes} on Discord) and notxxdog (\texttt{.notxxdog} on Discord).}

\begin{document}

\begin{abstract}
We study which integers $m$ satisfy $m\mid 2^{m-1}+3$, equivalently which $n$ satisfy
$n+1\mid 2^{n}+3$ (OEIS A245728). Only two are known---$61\cdot 228806497$ and
$67\cdot 44210291368986343$, each a product of two primes---and none lies below $10^{15}$. Our
central result ties the question to the prime factorisation of $m$: a squarefree
$m=p_{1}\cdots p_{\omega}$ is a solution if and only if, for every $i$, $-3$ is a power of $2$
modulo $p_{i}$ and $m\equiv k_{i}+1\pmod{d_{i}}$, where $d_{i}=\ord_{p_{i}}(2)$ and
$2^{k_{i}}\equiv-3\pmod{p_{i}}$. Equivalently, the largest prime factor of $m$ must divide
$2^{P-1}+3$ with $P$ the product of the remaining primes, so locating solutions reduces to
factoring a single, rapidly growing integer. These local conditions exclude a positive
proportion of primes as factors and cannot be sharpened by prime-power moduli; from them we show
that every solution is squarefree at each prime below $1093$, and that the only two-prime
solutions with smaller prime at most $499$ are the two above. Heuristics suggest that this set is infinite (\S\ref{sec:finite})---though extraordinarily
sparse, dominated by many-prime $m$ beyond any feasible factorisation, which is precisely why
only two-prime examples have been found. By contrast, $m\mid 2^{m-1}+1$ has no solutions, by a
short $2$-adic argument.
\end{abstract}

\maketitle
\tableofcontents

\section{Introduction}

For an integer $a$, write $S_{a}=\{\,n\ge 1:(n+1)\mid 2^{n}-a\,\}$. Setting $m=n+1$, membership
$n\in S_{a}$ is equivalent to
\begin{equation}\label{eq:main}
m\mid 2^{m-1}-a,\qquad m\ge 2.
\end{equation}
This note concerns $a=-1$ and $a=-3$. We prove $S_{-1}=\varnothing$ and assemble, both
structurally and computationally, what is presently known about $S_{-3}$.

\paragraph{Notation.}
For an odd prime $p$ with $p\nmid 2$, write $d_{p}=\ord_{p}(2)$, the multiplicative order of $2$
in $\F_{p}^{\times}$. For $a\in\Z$ with $\gcd(a,p)=1$ and $a\in\gen{2}\subseteq\F_{p}^{\times}$,
write $\log_{2}a$ for the smallest non-negative integer $k$ with $2^{k}\equiv a\pmod{p}$. We use
$\widetilde{\log_{2}}a$ analogously for the discrete logarithm in $(\Z/p^{2})^{\times}$. The
$2$-adic valuation of an integer $N$ is $v_{2}(N)$.

\section{The case \texorpdfstring{$a=-1$}{a=-1}}

\begin{theorem}\label{thm:minus1}
There is no integer $m>1$ with $m\mid 2^{m-1}+1$. Equivalently, $S_{-1}=\varnothing$.
\end{theorem}

\begin{proof}
Suppose such an $m$ exists. The integer $2^{m-1}+1$ is odd, so $m$ is odd. Fix any prime
$p\mid m$ and set $d=d_{p}$. From $p\mid 2^{m-1}+1$,
\[
2^{m-1}\equiv-1\pmod{p},\qquad 2^{2(m-1)}\equiv 1\pmod{p}.
\]
Hence $d\mid 2(m-1)$ but $d\nmid m-1$ (since $p$ is odd, $-1\not\equiv 1\pmod{p}$). It follows
that $d/\gcd(d,m-1)=2$, so $\gcd(d,m-1)=d/2$. Comparing $2$-adic valuations,
$v_{2}(d)-1=v_{2}(\gcd(d,m-1))\le v_{2}(m-1)$ on the one hand, while $d\nmid m-1$ together with
$d\mid 2(m-1)$ forces $v_{2}(d)>v_{2}(m-1)$ on the other. Therefore
\begin{equation}\label{eq:val}
v_{2}(d)=v_{2}(m-1)+1.
\end{equation}
Now $d\mid p-1$, so $v_{2}(p-1)\ge v_{2}(d)=v_{2}(m-1)+1$. Setting $\nu=v_{2}(m-1)+1$, we have
shown $p\equiv 1\pmod{2^{\nu}}$ for every prime $p\mid m$. Since $m$ is a product of primes in
$1+2^{\nu}\Z$, we get $m\equiv 1\pmod{2^{\nu}}$, i.e.\ $2^{\nu}\mid m-1$. But
$\nu=v_{2}(m-1)+1$, contradicting the definition of $v_{2}(m-1)$.
\end{proof}

\begin{remark}\label{rem:roots}
The proof gives slightly more: for any $m\ge 2$ and any odd $a$ with $\gcd(a,m)=1$ and
$a\equiv-1\pmod{p}$ for every $p\mid m$, the same $2$-adic argument applies. The obstruction is
genuinely about square roots of unity, and this is exactly why it does not transfer to $a=-3$:
$-3$ is not a root of unity modulo a typical prime, so no global multiplicative collapse occurs.
\end{remark}

\section{The case \texorpdfstring{$a=-3$}{a=-3}: elementary restrictions}

We turn to \eqref{eq:main} with $a=-3$, that is,
\begin{equation}\label{eq:star}\tag{$\star$}
m\mid 2^{m-1}+3.
\end{equation}

\begin{proposition}\label{prop:elem}
Let $m\ge 2$ satisfy \eqref{eq:star}. Then:
\begin{enumerate}[label=\textup{(\roman*)}]
\item $m$ is odd, equivalently $n=m-1$ is even;
\item $3\nmid m$;
\item $m$ is composite.
\end{enumerate}
\end{proposition}

\begin{proof}
(i) $2^{m-1}+3$ is odd, so $m$ is odd. (ii) If $3\mid m$ then $3\mid 2^{m-1}+3$, hence
$3\mid 2^{m-1}$, impossible. (iii) If $m$ is prime, Fermat's little theorem gives
$2^{m-1}\equiv 1\pmod m$, so $m\mid 4$, forcing $m\in\{2\}$, contradicting (i).
\end{proof}

\section{The known solutions}

The integers $n$ satisfying $(n+1)\mid 2^{n}+3$ form OEIS sequence A245728~\cite{oeis}. Only two
terms are known:
\[
m_{1}=n_{1}+1=13957196317=61\cdot 228806497,
\]
\[
m_{2}=n_{2}+1=2962089521722084981=67\cdot 44210291368986343,
\]
and no further $n$ exists with $n+1<10^{15}$. Both factorisations are products of two distinct
primes; in each case the smaller prime is small ($61$ or $67$) and the cofactor is large. Both
were verified directly to satisfy \eqref{eq:star}, and all four factors are prime. The second
solution exceeds $10^{18}$ and was found structurally rather than by enumeration; see
Section~\ref{sec:comp}.

\section{Admissibility: a sieve on prime factors}\label{sec:sieve}

The next lemma gives a strong necessary condition that every prime divisor of $m$ must satisfy;
it underlies all subsequent analysis.

\begin{lemma}[Local condition at a prime divisor]\label{lem:local}
Let $m\ge 5$ satisfy \eqref{eq:star}, let $p$ be a prime dividing $m$, and write $m=p^{a}t$ with
$\gcd(p,t)=1$. Then $p\notin\{2,3\}$, the element $-3$ lies in $\gen{2}\subseteq\F_{p}^{\times}$,
and, writing $d=d_{p}$ and $k=\log_{2}(-3)\bmod p$,
\begin{equation}\label{eq:loc}
t\equiv k+1\pmod{d}.
\end{equation}
\end{lemma}

\begin{proof}
By Proposition~\ref{prop:elem}, $p$ is odd and $p\ne 3$, so $\gcd(p,6)=1$. Since $d\mid p-1$,
$p\equiv 1\pmod d$, hence $p^{a}\equiv 1\pmod d$ and $m\equiv t\pmod d$. Therefore
$m-1\equiv t-1\pmod d$ and $2^{m-1}\equiv 2^{t-1}\pmod p$. Since $p\mid 2^{m-1}+3$, this gives
$2^{t-1}\equiv-3\pmod p$, which both shows $-3\in\gen{2}$ and yields \eqref{eq:loc}.
\end{proof}

\begin{definition}\label{def:adm}
A prime $p\notin\{2,3\}$ is \emph{admissible} if $-3\in\gen{2}\subseteq\F_{p}^{\times}$ and,
writing $d=d_{p}$ and $k=\log_{2}(-3)$:
\begin{enumerate}[label=\textup{(\Alph*)}]
\item if $d$ is even, then $k$ is even;
\item if $3\mid d$, then $k\not\equiv 2\pmod 3$.
\end{enumerate}
\end{definition}

\begin{theorem}\label{thm:adm}
Every prime divisor of any $m$ satisfying \eqref{eq:star} is admissible.
\end{theorem}

\begin{proof}
Fix $p\mid m$ and use the notation of Lemma~\ref{lem:local}, with $t=m/p^{v_{p}(m)}$. Since $m$
is odd, $t$ is odd; since $3\nmid m$, $3\nmid t$. Reduce \eqref{eq:loc} modulo $2$ and modulo
$3$. If $d$ is even, then $t\bmod d$ has the same parity as $t$, namely odd, so $k+1$ is odd,
i.e.\ $k$ is even. If $3\mid d$, then $t\bmod d$ has the same residue modulo $3$ as $t$, which is
$1$ or $2$ but not $0$, so $k+1\not\equiv 0\pmod 3$, i.e.\ $k\not\equiv 2\pmod 3$.
\end{proof}

\paragraph{Numerical effect of the sieve.}
Among the primes $p<200$ for which $-3\in\gen{2}$ in $\F_{p}^{\times}$ (a necessary condition
without which $p$ cannot divide any $m$), Theorem~\ref{thm:adm} eliminates substantially more
than half. Table~\ref{tab:sieve} records the data; admissible primes are marked $\checkmark$.
The admissible primes below $200$ are
\[
13,\ 19,\ 61,\ 67,\ 79,\ 103,\ 193,\ 199.
\]
The two known solutions use the admissible primes $61$ and $67$ together with admissible large
cofactors (one verifies directly that $228806497$ and $44210291368986343$ are admissible).
Condition (A) eliminates roughly half of the eligible primes (those with $d$ even and $k$ odd,
which by equidistribution is half), and condition (B) eliminates a further third of those with
$3\mid d$. Empirically, over primes below $2000$, about $56\%$ are $-3$-eligible and the sieve
removes about $61.5\%$ of those, leaving roughly $21.6\%$ of all primes admissible. Thus
admissible primes form a sparse but \emph{positive-proportion} subsequence (conjecturally of
positive density, by an Artin-type heuristic); in particular the supply of candidate small
factors is infinite, which already signals that no finite search can be conclusive.

\begin{table}[t]
\centering
\small
\begin{tabular}{rrrccc}
\toprule
$p$ & $d_{p}$ & $k$ & (A) $k$ even if $2\mid d$ & (B) $k\not\equiv 2\ (3)$ if $3\mid d$ & adm. \\
\midrule
5   & 4  & 1  & fail & ---  & $\times$ \\
7   & 3  & 2  & ---  & fail & $\times$ \\
11  & 10 & 3  & fail & ---  & $\times$ \\
13  & 12 & 10 & ok   & ok   & $\checkmark$ \\
19  & 18 & 4  & ok   & ok   & $\checkmark$ \\
29  & 28 & 19 & fail & ---  & $\times$ \\
37  & 36 & 8  & ok   & fail & $\times$ \\
53  & 52 & 43 & fail & ---  & $\times$ \\
59  & 58 & 21 & fail & ---  & $\times$ \\
61  & 60 & 36 & ok   & ok   & $\checkmark$ \\
67  & 66 & 6  & ok   & ok   & $\checkmark$ \\
79  & 39 & 10 & ---  & ok   & $\checkmark$ \\
83  & 82 & 31 & fail & ---  & $\times$ \\
97  & 48 & 43 & fail & ok   & $\times$ \\
101 & 100& 19 & fail & ---  & $\times$ \\
103 & 51 & 9  & ---  & ok   & $\checkmark$ \\
107 & 106& 17 & fail & ---  & $\times$ \\
131 & 130& 7  & fail & ---  & $\times$ \\
139 & 138& 110& ok   & fail & $\times$ \\
149 & 148& 13 & fail & ---  & $\times$ \\
163 & 162& 20 & ok   & fail & $\times$ \\
173 & 172& 113& fail & ---  & $\times$ \\
179 & 178& 19 & fail & ---  & $\times$ \\
181 & 180& 146& ok   & fail & $\times$ \\
193 & 96 & 90 & ok   & ok   & $\checkmark$ \\
197 & 196& 83 & fail & ---  & $\times$ \\
199 & 99 & 85 & ---  & ok   & $\checkmark$ \\
\bottomrule
\end{tabular}
\caption{Primes $p<200$ with $-3\in\gen{2}\pmod p$. Column (A) is vacuous when $d_{p}$ is odd;
column (B) is vacuous when $3\nmid d_{p}$.}
\label{tab:sieve}
\end{table}

\section{Interlocking congruences}\label{sec:interlock}

The local condition \eqref{eq:loc} globalises by the Chinese Remainder Theorem.

\begin{proposition}\label{prop:crt}
Let $m$ satisfy \eqref{eq:star} with prime factorisation $m=\prod_{i=1}^{r}p_{i}^{a_{i}}$. For
each $i$, set $d_{i}=d_{p_{i}}$ and $k_{i}=\log_{2}(-3)\bmod p_{i}$, and put $t_{i}=m/p_{i}^{a_{i}}$.
Then $t_{i}\equiv k_{i}+1\pmod{d_{i}}$ for every $i$, and consequently
\[
m\equiv p_{i}^{a_{i}}(k_{i}+1)\pmod{p_{i}^{a_{i}}d_{i}}
\]
for every $i$. By CRT, $m$ lies in a single residue class modulo
$\lcm_{i}\!\big(p_{i}^{a_{i}}d_{i}\big)$.
\end{proposition}

\begin{proof}
The first statement is Lemma~\ref{lem:local}. For the second, $m=p_{i}^{a_{i}}t_{i}$ and
$t_{i}\equiv k_{i}+1\pmod{d_{i}}$, so $m\equiv p_{i}^{a_{i}}(k_{i}+1)\pmod{p_{i}^{a_{i}}d_{i}}$.
\end{proof}

For both known solutions one verifies Proposition~\ref{prop:crt} explicitly:
\[
m_{1}=61\cdot 228806497:\quad 228806497 \bmod 60 = 37 = k_{61}+1,
\]
\[
m_{2}=67\cdot 44210291368986343:\quad 44210291368986343 \bmod 66 = 7 = k_{67}+1.
\]

\section{Two-prime structure}\label{sec:twoprime}

Both known $m$ are products of two distinct primes. Specialising Proposition~\ref{prop:crt}
gives a tight description of all such solutions, and the next remark shows that one half of the
description is automatic.

\begin{corollary}\label{cor:twoprime}
Suppose $m=pq$ with primes $p<q$, $p,q\ne 2,3$. Then $m$ satisfies \eqref{eq:star} if and only
if both
\begin{equation}\label{eq:pair}
q\equiv k_{p}+1\pmod{d_{p}},\qquad p\equiv k_{q}+1\pmod{d_{q}}
\end{equation}
hold, where $d_{p},d_{q}$ are the orders of $2$ and $k_{p},k_{q}$ the discrete logarithms of
$-3$ to base $2$ in $\F_{p}^{\times}$ and $\F_{q}^{\times}$ respectively. In particular $q$ is a
prime divisor of $2^{p-1}+3$ and $p$ is a prime divisor of $2^{q-1}+3$.
\end{corollary}

\begin{remark}[One-sidedness]\label{rem:onesided}
The two conditions in \eqref{eq:pair} are not independent once $q$ is drawn from the prime
factors of $N_{p}:=2^{p-1}+3$. Indeed, if $q\mid N_{p}$ then $2^{p-1}\equiv-3\pmod q$, and a
short computation with Fermat's little theorem shows
\[
2^{pq-1}\equiv (2^{p-1})^{q}\cdot 2^{q-1}\equiv(-3)^{q}\cdot 1\equiv -3\pmod q,
\]
so the congruence modulo $q$ holds automatically. Equivalently, the second condition
$p\equiv k_{q}+1\pmod{d_{q}}$ holds for \emph{every} prime factor $q$ of $N_{p}$. Hence
\[
\boxed{\,m=pq\ (p<q)\text{ solves }\eqref{eq:star}\iff q\mid 2^{p-1}+3\ \text{ and }\ q\equiv k_{p}+1\pmod{d_{p}}.\,}
\]
This makes the two-prime problem for a fixed smaller prime $p$ entirely a question about the
prime factorisation of the single integer $N_{p}$, together with one residue test modulo
$d_{p}$ ($<p$).
\end{remark}

This is the basis of an efficient algorithm: for each admissible small prime $p$, factor
$N_{p}=2^{p-1}+3$, and for each prime factor $q>p$ test whether $q\equiv k_{p}+1\pmod{d_{p}}$.

\section{Solutions with more prime factors}\label{sec:omega}

Corollary~\ref{cor:twoprime} is the case $\omega(m)=2$. The same reduction works for any number
of prime factors; it both clarifies the structure of hypothetical longer solutions and explains
why only two-prime ones are known. Throughout, $m=p_{1}p_{2}\cdots p_{\omega}$ is squarefree with
$p_{1}<p_{2}<\cdots<p_{\omega}$ distinct primes, $d_{i}=d_{p_{i}}$ and $k_{i}=\log_{2}(-3)\bmod
p_{i}$. (By Section~\ref{sec:sqfree} every solution is squarefree at all primes below $1093$, so
a repeated prime factor, if one ever occurs, must be $\ge 1093$; we return to this at the end.)

\begin{proposition}[Squarefree solutions of any length]\label{prop:omega}
Let $m=p_{1}\cdots p_{\omega}$ be squarefree with each $p_{i}\notin\{2,3\}$. Then
$m\mid 2^{m-1}+3$ if and only if, for every $i$,
\[
-3\in\gen{2}\subseteq\F_{p_{i}}^{\times}\qquad\text{and}\qquad m\equiv k_{i}+1\pmod{d_{i}},
\]
the latter being equivalent to $\displaystyle\prod_{j\ne i}p_{j}\equiv k_{i}+1\pmod{d_{i}}$. In
particular every $p_{i}$ is admissible.
\end{proposition}

\begin{proof}
Since $m$ is squarefree, the CRT gives $m\mid 2^{m-1}+3$ iff $p_{i}\mid 2^{m-1}+3$ for every
$i$. For fixed $i$, $2^{m-1}\equiv-3\pmod{p_{i}}$ holds iff $-3\in\gen{2}$ in $\F_{p_{i}}^{\times}$
and $m-1\equiv k_{i}\pmod{d_{i}}$. As $d_{i}\mid p_{i}-1$ we have $p_{i}\equiv 1\pmod{d_{i}}$,
hence $m\equiv\prod_{j\ne i}p_{j}\pmod{d_{i}}$, giving the reformulation. Conditions (A),(B) of
Definition~\ref{def:adm} follow as in Theorem~\ref{thm:adm}.
\end{proof}

Thus a squarefree solution is exactly a set of admissible primes whose product lands
simultaneously in the prescribed class $k_{i}+1$ modulo each $d_{i}$. By CRT this confines $m$ to
one residue class modulo $\lcm_{i}d_{i}$ (Proposition~\ref{prop:crt}); the difficulty is that the
unknown $m$ is forced to be the \emph{product} of the very primes defining the moduli. The
two-prime reduction of Remark~\ref{rem:onesided} generalises verbatim: the condition at the
largest prime is a divisibility, and the remaining ones are residue tests.

\begin{corollary}[Recursive criterion]\label{cor:omega}
Let $p_{1}<\cdots<p_{\omega}$ be admissible primes, put $P=p_{1}\cdots p_{\omega-1}=m/p_{\omega}$,
and write $N_{P}=2^{P-1}+3$. Then $m=Pp_{\omega}$ solves \eqref{eq:star} if and only if
\[
p_{\omega}\mid N_{P}\qquad\text{and}\qquad m\equiv k_{i}+1\pmod{d_{i}}\quad(1\le i\le\omega-1).
\]
The congruence $m\equiv k_{\omega}+1\pmod{d_{\omega}}$ is automatic once $p_{\omega}\mid N_{P}$.
\end{corollary}

\begin{proof}
The congruence at $p_{\omega}$ reads $P\equiv k_{\omega}+1\pmod{d_{\omega}}$ (since $m\equiv P
\pmod{d_{\omega}}$), i.e.\ $2^{P-1}\equiv-3\pmod{p_{\omega}}$, i.e.\ $p_{\omega}\mid N_{P}$; this
holds for every prime factor of $N_{P}$. The remaining $\omega-1$ congruences are
Proposition~\ref{prop:omega}.
\end{proof}

This is an algorithm: to find every $\omega$-prime solution extending a fixed tuple
$(p_{1},\dots,p_{\omega-1})$, factor the single integer $N_{P}=2^{P-1}+3$ and test each prime
factor $p_{\omega}>p_{\omega-1}$ against the $\omega-1$ residue conditions modulo
$d_{1},\dots,d_{\omega-1}$. Two points deserve emphasis. First, the smaller primes need
\emph{not} themselves form a solution: solutions do not nest, so one cannot grow a long solution
by extending a short one. Second, everything reduces to factoring $N_{P}$, and that is precisely
where the method fails.

\subsection{The computational barrier}
The integer $N_{P}=2^{P-1}+3$ has about $0.301\,P$ decimal digits, and $P=p_{1}\cdots
p_{\omega-1}$ is a product of admissible primes, hence grows super-exponentially with $\omega$.
Even the smallest admissible primes $13,19,61,67,\dots$, which minimise $P$, give the sizes of
Table~\ref{tab:growth}.

\begin{table}[h]
\centering
\small
\begin{tabular}{cccl}
\toprule
$\omega$ & minimal $P=p_{1}\cdots p_{\omega-1}$ & digits of $N_{P}$ & feasibility \\
\midrule
$2$ & $13$ & $4$ & routine (broad search) \\
$3$ & $13\cdot 19=247$ & $75$ & at the edge (one tuple) \\
$4$ & $13\cdot 19\cdot 61=15067$ & $\approx 4{,}500$ & infeasible \\
$5$ & $13\cdot 19\cdot 61\cdot 67=1{,}009{,}489$ & $\approx 3\times 10^{5}$ & infeasible \\
\bottomrule
\end{tabular}
\caption{Smallest integer $N_{P}=2^{P-1}+3$ that must be factored to search for an
$\omega$-prime solution.}
\label{tab:growth}
\end{table}

So two prime factors is the only regime open to a broad search, three is reachable for a single
tuple, and four or more is permanently beyond factorisation. Concretely, for three primes only
the smallest admissible pair $(p_{1},p_{2})=(13,19)$ keeps $N_{P}$ tractable:
\[
N_{247}=2^{246}+3=307\cdot 41047\cdot r,\qquad r\ \text{a 67-digit prime},
\]
and none of $13\cdot 19\cdot p_{3}$ with $p_{3}\in\{307,41047,r\}$ satisfies the two residue
conditions, so there is no three-prime solution with smaller pair $(13,19)$. Every other pair
already pushes $N_{P}$ past $200$ digits, so a complete three-prime search---let alone
$\omega\ge 4$---is not currently possible.

\subsection{How many solutions, and how large?}
Because the multi-prime regime is computationally invisible, the only guide is heuristic.
Modelling $2^{m-1}+3\bmod m$ as equidistributed makes a squarefree $m$ with all prime factors
admissible a solution with ``probability'' about $1/m$; when $d_{i}\approx p_{i}$ this agrees
with the structural estimate $\prod_{i}1/d_{i}$. Summing over candidates with exactly $\omega$
prime factors gives, for the number of $\omega$-prime solutions up to $X$,
\[
E_{\omega}(X)\ \approx\ \frac{1}{\omega!}\Bigl(\sum_{\substack{p\le X\\\text{admissible}}}
\frac{1}{p}\Bigr)^{\!\omega}\ \approx\ \frac{(\delta\,\log\log X)^{\omega}}{\omega!},
\]
where $\delta$ is the (logarithmic) density of admissible primes, empirically
$\delta\approx 0.15$. Hence
\[
\sum_{\omega\ge 2}E_{\omega}(X)\ \approx\ e^{\delta\log\log X}-1-\delta\log\log X
\ =\ (\log X)^{\delta}-1-\delta\log\log X\ \xrightarrow[X\to\infty]{}\ \infty.
\]
Three conclusions follow, of which the third is the point. (i)~The total diverges, so the
heuristic predicts \emph{infinitely many} solutions---but only as a minuscule power
$(\log X)^{\delta}$ of $\log X$, so they are astronomically sparse. (ii)~The estimate is of the
right order: at the scale of the known terms ($X\approx m_{2}$) it gives $\sum_{\omega\ge 2}
E_{\omega}$ of order $10^{-1}$, so finding two solutions is consistent with the model to within
its (considerable) crudeness; the structural refinement $\prod_{i}1/d_{i}$, which exceeds $1/m$
when some order $d_{i}$ is small, raises the estimate toward the observed count. (iii)~At every
reachable scale $\delta\log\log X<1$, so $E_{\omega}$ \emph{decreases} with $\omega$: two-prime
solutions are heuristically the most common, and a solution with three or more prime factors is
not expected until $X$ is of the order $\exp\!\bigl(\exp(c/\delta)\bigr)$---a number with
thousands of digits or more. The two known solutions are two-prime both because two-prime
solutions genuinely dominate at small heights and because, even were a longer solution to lie in
range, exhibiting it would require factoring an $N_{P}$ of thousands of digits. Longer solutions
are thus doubly hidden: too rare to meet in any searchable range, and walled off by
factorisation when they do occur.

\subsection{Non-squarefree \texorpdfstring{$m$}{m}}
If $p^{a}\,\|\,m$ with $a\ge 2$, Proposition~\ref{prop:omega} must be read modulo $p^{a}$: the
condition at $p$ becomes $2^{m-1}\equiv-3\pmod{p^{a}}$, governed by $\ord_{p^{a}}(2)$ and the
discrete logarithm of $-3$ there, in the manner of Lemma~\ref{lem:sqfree}. Corollary
\ref{cor:sqfree} forbids $a\ge 2$ for every prime $p<1093$, so a non-squarefree solution would
have to involve a repeated prime $\ge 1093$; none is known, and the heuristic above, dominated by
squarefree $m$, is unaffected.

\section{Squarefreeness at small primes}\label{sec:sqfree}

Both known solutions are squarefree. We now prove that this is forced at every non-Wieferich
prime, modulo a single discrete-logarithm coincidence that we then verify never occurs below
$1093$. Recall that $p$ is a \emph{Wieferich prime} (base $2$) if $2^{p-1}\equiv 1\pmod{p^{2}}$;
the only such primes below $4\times 10^{3}$ are $1093$ and $3511$. For $p$ non-Wieferich one has
$\ord_{p^{2}}(2)=p\,d_{p}$ (if $2^{d_{p}}\equiv 1\pmod{p^{2}}$ then $2^{p-1}\equiv 1\pmod{p^{2}}$,
since $d_{p}\mid p-1$; the contrapositive gives the order $p\,d_{p}$).

\begin{lemma}[Squarefreeness at a non-Wieferich prime]\label{lem:sqfree}
Let $m$ satisfy \eqref{eq:star} and let $p$ be a prime with $p^{2}\mid m$. Suppose $p$ is not a
Wieferich prime base $2$, so $D:=\ord_{p^{2}}(2)=p\,d_{p}$. Then $-3\in\gen{2}\subseteq
(\Z/p^{2})^{\times}$ and, writing $\widetilde{k}_{p}:=\widetilde{\log_{2}}(-3)$ for its discrete
logarithm modulo $p^{2}$,
\[
\widetilde{k}_{p}\equiv -1\pmod{p}.
\]
\end{lemma}

\begin{proof}
From $p^{2}\mid m\mid 2^{m-1}+3$ we get $2^{m-1}\equiv-3\pmod{p^{2}}$, so $-3\in\gen{2}$ modulo
$p^{2}$ and
\[
m-1\equiv\widetilde{k}_{p}\pmod{D}.
\]
Since $p$ is non-Wieferich, $D=p\,d_{p}$, so $p\mid D$ and we may reduce the previous congruence
modulo $p$:
\[
m-1\equiv\widetilde{k}_{p}\pmod{p}.
\]
But $p^{2}\mid m$ gives $m\equiv 0\pmod p$, hence $m-1\equiv-1\pmod p$. Therefore
$\widetilde{k}_{p}\equiv-1\pmod p$.
\end{proof}

\begin{remark}
The mechanism here differs from that of the sieve. Conditions (A),(B) of
Definition~\ref{def:adm} arise from the cofactor $t=m/p^{a}$ being odd and coprime to $3$, which
constrains $t$ only modulo $2$ and $3$; the classes modulo $4,9,\dots$ remain free
(Section~\ref{sec:nosharp}). In Lemma~\ref{lem:sqfree} the modulus is the prime $p$ itself, and
the order $D=p\,d_{p}$ ties the residue $m\bmod p$, which is $0$ since $p^{2}\mid m$, to
$\widetilde{k}_{p}\bmod p$. This is a genuine constraint rather than a free residue class.
\end{remark}

\begin{corollary}\label{cor:sqfree}
For every admissible prime $p<1093$ one has $\widetilde{k}_{p}\not\equiv-1\pmod p$. Consequently
no solution $m$ of \eqref{eq:star} is divisible by $p^{2}$ for any prime $p<1093$; that is, every
$m\in S_{-3}$ is squarefree at all primes below $1093$.
\end{corollary}

\begin{proof}
If $p^{2}\mid m$ then $p\mid m$, so $p$ is admissible by Theorem~\ref{thm:adm}, and every
admissible $p<1093$ is non-Wieferich. By Lemma~\ref{lem:sqfree} (and noting that if $-3\notin
\gen{2}$ modulo $p^{2}$ then $p^{2}\nmid m$ outright), divisibility by $p^{2}$ would force
$\widetilde{k}_{p}\equiv-1\pmod p$. A direct computation of the discrete logarithm
$\widetilde{k}_{p}$ for each of the $39$ admissible primes $p<1093$ shows
$\widetilde{k}_{p}\not\equiv-1\pmod p$ in every case, a contradiction.
\end{proof}

Table~\ref{tab:sqfree} records $\widetilde{k}_{p}\bmod p$ for the small admissible primes; the
value $-1\bmod p=p-1$ never appears.

\begin{table}[t]
\centering
\small
\begin{tabular}{rrrr}
\toprule
$p$ & $d_{p}$ & $\widetilde{k}_{p}\bmod p$ & $-1\bmod p$ \\
\midrule
13  & 12  & 7   & 12  \\
19  & 18  & 6   & 18  \\
61  & 60  & 19  & 60  \\
67  & 66  & 35  & 66  \\
79  & 39  & 7   & 78  \\
103 & 51  & 66  & 102 \\
193 & 96  & 7   & 192 \\
199 & 99  & 35  & 198 \\
211 & 210 & 186 & 210 \\
\bottomrule
\end{tabular}
\caption{For small admissible primes, the reduced discrete logarithm
$\widetilde{k}_{p}=\widetilde{\log_{2}}(-3)\bmod p^{2}$ taken modulo $p$, against $-1\bmod p$.
Equality (which would permit $p^{2}\mid m$) never occurs; verified for all admissible $p<1093$.}
\label{tab:sqfree}
\end{table}

A full proof of squarefreeness would also require a uniform treatment of large primes and a
separate analysis of the Wieferich primes $1093$ and $3511$; the present argument settles only
the primes below $1093$.

\section{No prime-power sharpening of the sieve}\label{sec:nosharp}

It is natural to ask whether replacing the moduli $2,3$ in conditions (A),(B) by higher prime
powers $4,8,\dots$ or $9,27,\dots$ yields additional exclusions. It does not.

\begin{proposition}\label{prop:nosharp}
The only constraints on $t=m/p^{a}$ available in isolation are $t$ odd and $3\nmid t$. Reducing
the local congruence \eqref{eq:loc} modulo any power $2^{e}$ or $3^{f}$ produces no exclusion
beyond conditions \textup{(A)} and \textup{(B)}.
\end{proposition}

\begin{proof}
Suppose $2^{e}\mid d$. Then \eqref{eq:loc} forces $t\equiv k+1\pmod{2^{e}}$. The constraint
$t$ odd permits exactly the odd residues modulo $2^{e}$, and the forced value $k+1$ fails to be
odd precisely when $k+1$ is even, i.e.\ when $k$ is odd. This is condition (A) (which already
applies as soon as $2\mid d$). Higher $e$ add nothing, since oddness is a condition modulo $2$
only. Similarly, if $3^{f}\mid d$ then $t\equiv k+1\pmod{3^{f}}$, and $3\nmid t$ permits exactly
the residues coprime to $3$; the forced value $k+1$ fails to be coprime to $3$ precisely when
$3\mid k+1$, i.e.\ $k\equiv 2\pmod 3$, which is condition (B). Higher $f$ add nothing.
\end{proof}

We verified computationally that no admissible prime below $3000$ is removed by any mod-$4$,
mod-$8$, mod-$9$, or mod-$27$ refinement, consistent with Proposition~\ref{prop:nosharp}. Any
genuine sharpening of the sieve must therefore come from \emph{interactions between distinct
prime factors}, i.e.\ from the interlocking congruences of Proposition~\ref{prop:crt}, rather
than from a single local modulus.

\section{Computational evidence}\label{sec:comp}

\subsection{Direct search}
The bound $m\le 10^{15}$ is the range exhaustively verified in OEIS A245728~\cite{oeis}; the only
solution there is $m_{1}$. An independent brute-force test over all odd $m\le 10^{7}$ with
$3\nmid m$ reproduces no solution in that range, agreeing with the database.

\subsection{Two-prime search via factorisation of \texorpdfstring{$2^{p-1}+3$}{2^(p-1)+3}}
By Remark~\ref{rem:onesided}, for each admissible smaller prime $p$ it suffices to factor
$N_{p}=2^{p-1}+3$ and test the prime factors $q>p$ against $q\equiv k_{p}+1\pmod{d_{p}}$. We
carried this out for every admissible $p\le 499$:
\begin{itemize}[nosep]
\item For $p\in\{13,19,61,67,79,103,193,199,211,271,307,373,463,499\}$ the number $N_{p}$ is
fully factored (by trial division, Pollard $\rho$, $p-1$, and ECM). The complete list of
two-prime solutions with smaller prime in this set is exactly $\{m_{1},m_{2}\}$, arising at
$p=61$ and $p=67$.
\item For $p\in\{367,439,487\}$ a composite cofactor of $109$, $132$, and $147$ decimal digits
respectively remains unfactored; in each case all prime factors found so far yield no solution,
and (each residual cofactor having been checked to be composite, not a single large prime) the
possibility of a further two-prime $m$ involving these $p$ is left open.
\end{itemize}

\begin{table}[t]
\centering
\small
\begin{tabular}{lll}
\toprule
smaller prime $p$ & $N_{p}=2^{p-1}+3$ status & solution \\
\midrule
$13,19$ & full & none \\
$61$ & full & $m_{1}=61\cdot 228806497$ \\
$67$ & full & $m_{2}=67\cdot 44210291368986343$ \\
$79,103,193,199,211,271,307,373,463,499$ & full & none \\
$367,439,487$ & partial ($109$--$147$-digit cofactor) & open \\
\bottomrule
\end{tabular}
\caption{Two-prime search over admissible smaller primes $p\le 499$. By Theorem~\ref{thm:adm}
only admissible primes can occur as a smaller factor, so non-admissible primes require no
factorisation.}
\label{tab:search}
\end{table}

A search for solutions with three or more prime factors is governed by the recursive criterion
of Section~\ref{sec:omega}; only the single tuple $(13,19)$ is within reach, and it yields no
three-prime solution.

\section{Is \texorpdfstring{$S_{-3}$}{S(-3)} finite?}\label{sec:finite}

This is the heart of ``fully solving'' the problem, and it is open. The available heuristics---a
structural estimate for two-prime solutions and a global density argument---point the same way,
toward infinitude; we give both and then say what a proof would require.

\paragraph{The two-prime estimate.}
For admissible $p$, the number $N_{p}=2^{p-1}+3$ has about $\log\log N_{p}\asymp\log p$ distinct
prime factors, and a given one lies in the class required modulo $d_{p}$ with ``probability''
about $1/d_{p}$. The expected number of two-prime solutions with smaller prime below $X$ is thus
of order $\sum_{p<X,\ \mathrm{adm}}(\log p)/d_{p}$, and this \emph{diverges}. Since $d_{p}\mid p-1$
we have $d_{p}<p$, hence $(\log p)/d_{p}>(\log p)/p$; assuming, as the data strongly support, that
the admissible primes have positive density $\delta$, Mertens' first theorem gives
\[
\sum_{\substack{p<X\\\text{admissible}}}\frac{\log p}{d_{p}}\ \ge\
\sum_{\substack{p<X\\\text{admissible}}}\frac{\log p}{p}\ \sim\ \delta\log X\ \xrightarrow[X\to\infty]{}\ \infty .
\]
The factor $\log p$ in the numerator is what fixes the rate at $\log X$: the bare sum $\sum
1/d_{p}$, lower-bounded by $\sum 1/p$, diverges only at the slower rate $\log\log X$ (numerically
$\sum_{p<X,\ \mathrm{adm}}1/d_{p}$ creeps from about $0.24$ at $X=200$ to $0.53$ at $X=5\cdot
10^{4}$). So the two-prime heuristic predicts \emph{infinitely many} two-prime solutions, not
finiteness. It is moreover well calibrated: $\sum_{p<X,\ \mathrm{adm}}(\log p)/d_{p}$ is about
$1.7$ at $X=2000$ and $2.4$ at $X=10^{4}$, consistent with the two known solutions, whose smaller
primes are $61$ and $67$.

\paragraph{A global density heuristic.}
Model $2^{m-1}+3\bmod m$ as roughly equidistributed; then $m$ is a solution with ``probability''
about $1/m$, and $\sum_{m\le X}1/m\sim\log X\to\infty$. A divergent expectation predicts
\emph{infinitely many} solutions. This is precisely the (correct) heuristic for the closely
analogous Fermat-pseudoprime problem $m\mid 2^{m-1}-2$, for which infinitely many solutions are
known (Erd\H{o}s). The inhomogeneous shift ($-3$ in place of $-1$) changes the combinatorics of
how prime factors glue together but not the divergence of the expectation.

Both estimates thus point the same way. The two-prime sum already diverges, and the global
heuristic is dominated by solutions with \emph{many} prime factors, whose expected distribution
(Section~\ref{sec:omega}) also diverges; together they favour an infinite $S_{-3}$ whose elements
beyond $m_{2}$ are many-prime integers far beyond any search. On this view the two known terms
being two-prime reflects the limits of factorisation, not the true distribution of solutions.

A proof of \emph{finiteness} would need an upper bound, uniform in $p$, on the number of prime
factors of $2^{p-1}+3$ lying in prescribed arithmetic progressions, which is well beyond current
analytic number theory. A proof of \emph{infinitude} would need a Carmichael-style construction
producing solutions to order; the homogeneous case $a=-1$ admits the relevant root-of-unity
structure (and is settled in the opposite direction, $S_{-1}=\varnothing$), but for $a=-3$ the
varying targets $k_{q}$ obstruct the standard construction and none is known. Both directions are
presently out of reach.

\section{Open questions}

\begin{enumerate}[label=\arabic*.]
\item \emph{Is $S_{-3}$ finite?} This remains open; the heuristics of Section~\ref{sec:finite}
point toward an infinite, though computationally unreachable, solution set.
\item \emph{Can the admissibility sieve be sharpened locally?} No: Proposition~\ref{prop:nosharp}
shows that prime-power moduli yield nothing further, so any improvement must be global.
\item \emph{Is every $m\in S_{-3}$ squarefree?} This is open in general;
Corollary~\ref{cor:sqfree} settles it at every prime below $1093$, and a uniform treatment of
large primes, together with a separate analysis of the Wieferich primes $1093,3511$, would be
needed for the full statement.
\item \emph{The same questions for $a\in\{-5,-7,\dots\}$.} The proof of Theorem~\ref{thm:minus1}
uses critically that $-1$ is a root of unity; for general $a$ the $2$-adic argument fails and the
qualitative behaviour of $S_{a}$ is unclear.
\end{enumerate}

\appendix

\section{Computational programs}\label{app:programs}
The programs used for the computations reported here---the direct search, the two-prime
factoring search via $N_{p}=2^{p-1}+3$, and the supporting order and discrete-logarithm
routines---are available at
\begin{center}
\url{https://github.com/RobinCodes/Projects/tree/main/Divisibility%20problem}.
\end{center}
The factorisations of Table~\ref{tab:search}, the discrete-logarithm verification underlying
Corollary~\ref{cor:sqfree}, and the sieve data of Table~\ref{tab:sieve} were carried out with
\texttt{sympy}.

\section{Note on the use of artificial intelligence}\label{app:ai}
This work was prepared with substantial assistance from an artificial-intelligence system, and
most of the results presented here---the admissibility sieve, the interlocking-congruence
description, the two-prime reduction, the squarefreeness analysis, and the computations of
Section~\ref{sec:comp}---were created with such assistance. We do not specify the particular
system used; readers seeking further detail about the methods are invited to contact the authors:
Robin Rovenszky (\texttt{robincodes} on Discord) and notxxdog (\texttt{.notxxdog} on Discord).

No result here is accepted on the basis of machine assertion. Every definition, lemma, and proof
is stated so as to be checkable by hand, and every computational claim is reproducible by
re-running the programs of Appendix~\ref{app:programs}.

\begin{thebibliography}{9}
\bibitem{oeis} OEIS Foundation Inc.\ (2024), Entry A245728 in \emph{The On-Line Encyclopedia of
Integer Sequences}, \url{https://oeis.org/A245728}.
\end{thebibliography}

\end{document}
